% Level 3, Week 4 answers -- Searching and How Fast Is Fast?
\answersweek{4}{Searching and How Fast Is Fast?}{Answers for workbook chapter 4 and homework 4.}

\section{Workbook exercises}
\begin{answer}[Your turn: trace it before you code it]
\begin{tabular}{@{}cccccl@{}}
\toprule
step & low & high & mid & items[mid] & so\ldots \\ \midrule
1 & 0 & 7 & 3 & 21 & too small: low = 4 \\
2 & 4 & 7 & 5 & 37 & too small: low = 6 \\
3 & 6 & 7 & 6 & 44 & too small: low = 7 \\
4 & 7 & 7 & 7 & 52 & found at index 7 \\
\bottomrule
\end{tabular}\\[4pt]
4 steps with binary search; linear search would take 8. Row 4 is the important one: \texttt{low} and \texttt{high} are equal and there is still one item to check (that is why the loop says \texttt{<=}).
\end{answer}

\subsection{Exercise 4.1: Search Race}
\pyfile{../exercises/week04/ex-search-race/solution.py}
\begin{answer}[Exercise 4.1: Search Race]
\textbf{Prediction:} 10000 items need 14 binary steps (the program prints 10, 100, 1000, 10000 linear steps and 4, 7, 10, 14 binary steps). The halving ladder is 10000, 5000, 2500, 1250, 625, 312, 156, 78, 39, 19, 9, 4, 2, 1: 14 numbers. A prediction of 13 from the \enquote{about 3 more per times 10} rule is a good estimate; praise the reasoning.\\
\textbf{What to look for:} \texttt{steps} goes up once per loop, before the comparison; \texttt{//} for \texttt{mid}; \texttt{mid + 1} and \texttt{mid - 1}.
A version that counts two comparisons per loop gives bigger step numbers and may fail the \enquote{10 steps for 1000 items} check; explain that we count one \enquote{look} per loop.
\end{answer}

\subsection{Exercise 4.2: Guess-the-Number Bot}
\pyfile{../exercises/week04/ex-guess-bot/solution.py}
\begin{answer}[Exercise 4.2: Guess-the-Number Bot]
\textbf{What to look for:} \texttt{next\_guess} returns \texttt{(low + high) // 2} (rounding up is also accepted); the bot never needs more than 7 guesses for 1 to 100. Typical games take 5 to 7 guesses; secret 50 takes 1.
If a student answers wrongly on purpose, \texttt{low} passes \texttt{high} and the bot says the answers do not add up: a nice discussion about input you cannot trust.
\end{answer}

\section{Homework}
\subsection{Core: Growth Table}
\pyfile{../exercises/week04/hw-growth-table/solution.py}
\begin{answer}[Core: Growth Table]
\begin{outputcode}
   items |  linear |  binary
      10 |      10 |       4
     100 |     100 |       7
    1000 |    1000 |      10
 1000000 | 1000000 |      20
\end{outputcode}
Predictions do not need to be right: what matters is that the student noticed the binary column grows by about 3 for each \enquote{times 10}.
\end{answer}

\begin{answer}[Challenge: Trace it by hand]
\begin{tabular}{@{}cccccl@{}}
\toprule
step & low & high & mid & items[mid] & so\ldots \\ \midrule
1 & 0 & 6 & 3 & 7 & too small: low = 4 \\
2 & 4 & 6 & 5 & 11 & too big: high = 4 \\
3 & 4 & 4 & 4 & 9 & found at index 4 \\
\bottomrule
\end{tabular}\\[4pt]
3 steps with binary search; linear search would take 5.
\end{answer}

\begin{answer}[Creative: Draw the race]
Any chart where the pairs of bars are labeled with the real numbers. Good ideas: a broken bar with a zig-zag, a different scale for each pair (clearly labeled), or bars for binary only next to written linear numbers. Discuss: a chart can mislead if the scale is hidden.
\end{answer}
